% Part: computability
% Chapter: recursive-functions
% Section: bounded-minimization

\documentclass[../../../include/open-logic-section]{subfiles}

\begin{document}

\olfileid{cmp}{rec}{bmi}
\olsection{محدوده کمينه‌موندنه}

\begin{explain}
ډېر وخت دا ګټوره وي چې يوه تابع د هغه تر ټولو کوچني عدد په توګه تعريف کړو چې کوم خاصيت يا اړيکه~$P$ پوره کوي۔ که~$P$ پرېکړه کېدونکې وي، دا تابع يوازې په دې ډول محاسبه کولے شو چې ټول ممکن عددونه~$0$، $1$، $2$، \dots يو په بل پسې وازمېيو، تر هغه چې تر ټولو کوچنے عدد ومومو چې~$P$ پوره کوي۔ دا ډول بې‌پولې پلټنه مو د بنسټيزو بازګشتي تابعو له چوکاټه باسي۔ خو که يوازې په هغه تر ټولو کوچني عدد کښې دلچسپي ولرو چې
\emph{له کوم خپلواکه ورکړل شوي حد څخه کوچنے وي}، نو بنسټيز بازګشتي پاتې کېږو۔ په بله او لږ عمومي وينا، فرض کړئ يوه بنسټيزه بازګشتي اړيکه~$R(x,z)$ لرو۔ هغه تابع په پام کښې ونيسئ چې $x$ او~$y$ هغه تر ټولو کوچني~$z < y$ ته لېږي چې~$R(x, z)$ پوره کوي۔ دا تابع هم د~$R(x, 0)$، $R(x, 1)$، \dots، $R(x, y-1)$ په ازمېيلو محاسبه کېدے شي۔ خو دا ولې بنسټيزه بازګشتي ده؟
\end{explain}

\begin{prop}
که~$R(\vec x, z)$ بنسټيزه بازګشتي وي، نو تابع~$m_R(\vec{x}, y)$ هم داسې ده؛ دا تابع، که داسې عدد وي، له~$y$ څخه تر ټولو کوچنے~$z$ راګرځوي چې~$R(\vec x, z)$ پوره کړي، او که داسې عدد نۀ وي نو~$y$ راګرځوي۔ تابع~$m_R$ به داسې ليکو:
\[
\bmin{z < y}{R(\vec{x}, z)},
\]
\end{prop}

\begin{proof}
داسې هېڅ~$z < 0$ نشته چې~$R(\vec{x}, z)$ پوره کړي، ځکه په اصل کښې هېڅ~$z < 0$ نشته۔ نو~$m_R(\vec x, 0) = 0$۔

کله چې حد د~$y + 1$ په بڼه وي، درې حالتونه لرو:
\begin{enumerate}
\item داسې~$z < y$ شته چې~$R(\vec{x}, z)$ پوره کوي؛ په دې حالت کښې $m_R(\vec{x}, y+1) = m_R(\vec{x}, y)$۔
\item داسې~$z<y$ نشته، خو~$R(\vec{x}, y)$ پوره کېږي؛ بيا $m_R(\vec{x}, y+1) = y$۔
\item داسې~$z < y+1$ نشته چې~$R(\vec{x}, z)$ پوره کړي؛ بيا $m_R(\vec{x}, y+1) = y+1$۔
\end{enumerate}
\textbf{د ژباړې د نسخې سمون (OLCMP-007):} سرچينې په درېيم حالت کښې د تابع د پاراميټرونو غلط وکټور ليکلی و، حال دا چې په ټول تعريف کښې هماغه لومړنی وکټور ثابت پاتې کېږي۔ دلته آرګومېنټ بېرته هماغه ثابت وکټور ته سم شوے دے۔
نو د~$m_R(\vec x, 0)$ پېلنی حالت په بنسټيز بازګښت د لاندې بشپړ تعريف په ترڅ کښې ورکولے شو:
\begin{align*}
m_R(\vec{x}, 0) & = 0\\
m_R(\vec{x}, y+1) & =
\begin{cases}
m_R(\vec{x}, y) & \text{که $m_R(\vec x, y) \neq y$ وي}\\
y & \text{که $m_R(\vec x, y) = y$ او $R(\vec{x}, y)$ وي}\\
y+1 & \text{که نۀ۔}
\end{cases}
\end{align*}
پام وکړئ چې داسې~$z<y$ شته چې~$R(\vec x, z)$ پوره کړي، هله او يوازې هله چې~$m_R(\vec x, y) \neq y$ وي۔
\end{proof}

\begin{prob}
فرض کړئ~$R(\vec x, z)$ بنسټيزه بازګشتي ده۔ تابع~$m'_R(\vec{x}, y)$ چې، که داسې عدد وي، له~$y$ څخه تر ټولو کوچنے~$z$ راګرځوي چې $R(\vec{x}, z)$ پوره کړي، او که داسې عدد نۀ وي نو~$0$ راګرځوي، له~$\Char{R}$ څخه په بنسټيز بازګښت تعريف کړئ۔
\end{prob}

\end{document}
